
Reinforcement Learning from Human Feedback (RLHF) and Direct Preference Optimization (DPO) represent two dominant paradigms for aligning language models with human preferences. These methods take fundamentally different optimization paths: RLHF trains an explicit reward model that learns continuous reward predictions, while DPO directly optimizes from preference pairs without an intermediate reward signal. This study tests whether these mechanistic differences produce distinct alignment signatures under controlled conditions using Llama-2-7B and Anthropic's HH-RLHF dataset.

A five-hypothesis experimental design was employed. The mechanistic hypotheses were validated: RLHF produces smooth, continuous reward predictions spanning a 0.83-unit range, while DPO preserves sharper preference boundaries with 65\% higher margin variance (sharpness ratio = 1.65). However, these training dynamics did not translate to distinct downstream effects at 7B scale with LoRA fine-tuning. Models trained with different methods did not cluster into separate behavioral attractors; cross-method similarity exceeded within-method similarity (clustering gap = -0.016). Standard alignment benchmarks (TruthfulQA, HHH) showed similar performance profiles (max |Cohen's d| = 0.194, profile correlation = 0.978).

The central finding is characterized as ''different dynamics, similar destinations'': verified mechanistic differences do not manifest as distinct alignment signatures on existing evaluation infrastructure at this scale. This has implications for method selection and suggests limitations in current benchmark sensitivity.

